\chapter{实验设置与结果分析}
\label{chap:exp}

本章给出实验配置、训练过程、测试集结果、消融实验分析以及端到端推理验证。

\section{实验设置}

\subsection{模型配置}

模型配置如表\ref{tab:model_config}所示。

\begin{table}[H]
\centering
\caption{模型配置}
\label{tab:model_config}
\begin{tabular}{ll}
\toprule
组件 & 配置 \\
\midrule
视频输入维度 & 768 \\
音频输入维度 & 768 \\
生理输入维度 & 256 \\
用药输入维度 & 128 \\
融合隐藏维度 & 512 \\
注意力头数 & 8 \\
Cross-Attention层数 & 4 \\
Dropout & 0.2（融合）/ 0.3（分类头） \\
分类头 & $512 \to 256 \to 128 \to 3$ \\
\bottomrule
\end{tabular}
\end{table}

\subsection{训练配置}

训练配置如表\ref{tab:train_config}所示。

\begin{table}[H]
\centering
\caption{训练配置}
\label{tab:train_config}
\begin{tabular}{ll}
\toprule
参数 & 值 \\
\midrule
优化器 & AdamW (lr=$10^{-4}$, weight\_decay=0.01) \\
学习率调度 & 5 epoch warmup + cosine annealing (min\_lr=$10^{-6}$) \\
Batch Size & 32 \\
最大Epochs & 30 \\
早停 & patience=10 \\
损失函数 & CrossEntropyLoss (label\_smoothing=0.1) \\
类别权重 & [1, 2, 3]（低/中/高） \\
梯度裁剪 & 1.0 \\
设备 & CPU \\
\bottomrule
\end{tabular}
\end{table}

\subsection{评估指标}

采用以下评估指标：准确率（Accuracy）、精确率/召回率/F1（Macro \& Weighted）、AUC（Macro, one-vs-rest）、混淆矩阵、各类别P/R/F1。

\subsection{消融实验设计}

消融实验分三组：

\begin{itemize}
    \item \textbf{A组 - 融合策略对比}：A1 Full（Cross-Modal Attention）vs A2 Lite（MLP拼接），对比跨模态注意力与简单拼接的效果
    \item \textbf{B组 - 单模态屏蔽}：分别将video/audio/health/medication特征置零，观察准确率下降幅度
    \item \textbf{C组 - 模态重要性量化}：$C_{importance} = A1_{acc} - B_{mask\_acc}$，下降越多说明该模态越重要
\end{itemize}

\section{训练过程}

Full模型在15个epoch后早停（patience=10无提升）。关键epoch的训练指标如表\ref{tab:training}所示。

\begin{table}[H]
\centering
\caption{训练过程关键指标}
\label{tab:training}
\begin{tabular}{ccccc}
\toprule
Epoch & Train Loss & Train Acc & Val Loss & Val Acc \\
\midrule
1 & 0.9545 & 40.3\% & 0.7032 & 56.7\% \\
3 & 0.4546 & 94.2\% & 0.3025 & 97.8\% \\
5 & 0.3104 & 99.1\% & 0.2819 & \textbf{100\%} \\
10 & 0.3221 & 98.4\% & 0.3386 & 100\% \\
15 & 0.2873 & 100\% & 0.2741 & 100\% \\
\bottomrule
\end{tabular}
\end{table}

模型在第5个epoch即达到验证集100\%准确率，此后验证损失持续微降（0.2819$\to$0.2741），训练损失稳定在0.28--0.31区间。训练与验证损失差距较小，表明模型收敛良好，无明显过拟合趋势。

\section{测试集结果}

Full模型在90样本独立测试集上的表现如表\ref{tab:test_results}所示。

\begin{table}[H]
\centering
\caption{测试集评估结果}
\label{tab:test_results}
\begin{tabular}{lc}
\toprule
指标 & 值 \\
\midrule
\textbf{准确率} & \textbf{1.0000 (100\%)} \\
精确率(Macro) & 1.0000 \\
召回率(Macro) & 1.0000 \\
F1(Macro) & 1.0000 \\
AUC(Macro) & 1.0000 \\
\bottomrule
\end{tabular}
\end{table}

混淆矩阵如表\ref{tab:confusion}所示，三个风险类别均达到完美的分类，无任何误分类。

\begin{table}[H]
\centering
\caption{测试集混淆矩阵}
\label{tab:confusion}
\begin{tabular}{lccc}
\toprule
 & 预测低 & 预测中 & 预测高 \\
\midrule
实际低 & 29 & 0 & 0 \\
实际中 & 0 & 28 & 0 \\
实际高 & 0 & 0 & 33 \\
\bottomrule
\end{tabular}
\end{table}

\section{消融实验}

\subsection{A组：融合策略对比}

A组对比Cross-Modal Attention（Full）与MLP拼接（Lite）两种融合策略，结果如表\ref{tab:ablation_a}所示。

\begin{table}[H]
\centering
\caption{A组：融合策略对比}
\label{tab:ablation_a}
\begin{tabular}{lrrl}
\toprule
模型 & 参数量 & 准确率 & 混淆矩阵 \\
\midrule
A1 Full (Cross-Attn) & 8,511,495 & \textbf{100.00\%} & [[29,0,0],[0,28,0],[0,0,33]] \\
A2 Lite (MLP Concat) & 3,740,163 & 94.44\% & [[29,0,0],[0,24,4],[0,1,32]] \\
\bottomrule
\end{tabular}
\end{table}

Full模型（Cross-Attention）准确率比Lite（MLP拼接）高5.56个百分点。Lite模型主要在中风险$\to$高风险方向有4个误判、高风险$\to$中风险1个误判。这表明Cross-Modal Attention的跨模态交互能力确实优于简单特征拼接——通过注意力机制，模型能够学习``哪些模态组合对当前判断更相关''，从而更准确地区分中风险与高风险。

\subsection{B组：单模态屏蔽}

B组分别屏蔽每个模态（将特征置零），观察准确率下降，结果如表\ref{tab:ablation_b}所示。

\begin{table}[H]
\centering
\caption{B组：单模态屏蔽实验}
\label{tab:ablation_b}
\begin{tabular}{lcccl}
\toprule
屏蔽模态 & 准确率 & F1-Macro & 准确率下降 & 主要误判 \\
\midrule
基线(无屏蔽) & 100.00\% & 1.0000 & --- & 无 \\
屏蔽 health & 67.78\% & 0.6250 & \textbf{-32.22\%} & 高$\to$中: 29个 \\
屏蔽 medication & 97.78\% & 0.9780 & -2.22\% & 中$\to$高: 1, 高$\to$中: 1 \\
屏蔽 video & 98.89\% & 0.9890 & -1.11\% & 中$\to$高: 1 \\
屏蔽 audio & 98.89\% & 0.9890 & -1.11\% & 中$\to$高: 1 \\
\bottomrule
\end{tabular}
\end{table}

\subsection{C组：模态重要性量化}

C组根据B组结果量化各模态重要性，排名如表\ref{tab:ablation_c}所示。

\begin{table}[H]
\centering
\caption{C组：模态重要性排名}
\label{tab:ablation_c}
\begin{tabular}{clll}
\toprule
排名 & 模态 & 重要性(准确率下降) & 分析 \\
\midrule
1 & health & -32.22\% & 最重要。生理特征按风险等级有明确数值差异 \\
2 & medication & -2.22\% & 第二。用药依从性提供辅助风险线索 \\
3 & video & -1.11\% & 视觉运动特征有贡献 \\
3 & audio & -1.11\% & 声学特征有贡献 \\
\bottomrule
\end{tabular}
\end{table}

\textbf{关键发现}：四模态全部参与决策（均有非零下降）。这与迭代2（视频模态贡献为0\%，因特征为零向量）形成鲜明对比，验证了真实视频特征提取的有效性。

屏蔽health后，33个高风险样本中有29个被误判为中风险。这是因为高风险的生理异常信号（心率严重异常、血氧$<$90\%、血压危急）丢失后，模型无法区分中风险与高风险——用药和视频/音频特征在中高风险间的区分度不如生理数据明显。这证明生理数据是区分中高风险的关键模态。

\section{端到端推理验证}

使用ESC-50真实音频在应用中进行端到端测试（视频使用默认特征），每个风险等级测试2个音频文件，结果如表\ref{tab:e2e}所示。

\begin{table}[H]
\centering
\caption{端到端推理验证结果}
\label{tab:e2e}
\begin{tabular}{lllc}
\toprule
风险等级 & 音频类别 & 预测结果 & 置信度 \\
\midrule
低风险 & dog & \checkmark 低风险 & 91.0\% \\
低风险 & dog & \checkmark 低风险 & 90.6\% \\
中风险 & coughing & \checkmark 中风险 & 95.3\% \\
中风险 & coughing & \checkmark 中风险 & 95.3\% \\
高风险 & crying\_baby & \checkmark 高风险 & 97.2\% \\
高风险 & glass\_breaking & \checkmark 高风险 & 97.4\% \\
\bottomrule
\end{tabular}
\end{table}

\textbf{结果：6/6正确（100\%）}，且置信度随风险等级递增（91\%$\to$95\%$\to$97\%），符合预期——高风险场景的声学特征（哭声、玻璃破碎声）区分度更高，模型更有信心。

\section{与迭代2对比}

本工作（迭代3）与迭代2的对比如表\ref{tab:compare}所示。

\begin{table}[H]
\centering
\caption{迭代2与迭代3对比}
\label{tab:compare}
\begin{tabular}{lll}
\toprule
维度 & 迭代2 & 迭代3（本工作） \\
\midrule
视频数据 & 零向量（URFD未下载） & Pexels 112视频（真实骨架） \\
视频特征非零率 & 7.6\% & 100\% \\
数据规模 & 380样本 & 900样本（三类均衡） \\
Full准确率 & 94.74\% & 100.00\% \\
视频模态贡献 & 0\% & -1.11\% \\
参与决策模态数 & 2/4 & 4/4 \\
MediaPipe API & mp.solutions（已废弃） & Tasks API (0.10.35) \\
\bottomrule
\end{tabular}
\end{table}

引入真实视频特征后，模型准确率提升5.26个百分点，视频模态从无贡献变为有贡献，四模态全部参与决策。这证明了真实数据特征对多模态融合模型性能的关键作用。
